\section{Commercial Bayesian and hidden-state engines: what a posterior can justify}
\label{exec22:state-scope}

A servicing decision may depend on a credit state that the bank cannot directly
observe. A missed payment, account balance and contact history provide evidence
about that state, but the proposed action changes a mechanism rather than merely
revealing another measurement. A useful commercial comparison must therefore
separate four questions: what state is plausible given information available now;
what parameters are plausible within a specified model; which model restrictions
the observations identify; and what happens when the bank changes an action.
An accurate answer to the first two questions need not answer the last two.

The relevant commercial landscape includes graphical-model engines, statistical
estimators and general numerical environments. Three complete Stata command
manuals add detail to the previously inspected GeNIe/SMILE, Netica, BayesiaLab and
Mplus profiles. These five supported families allow a substantive comparison,
while SAS, EViews, HUGIN, Bayes Server and commercial numerical environments
remain candidates requiring further matched technical evidence. The comparison
retains the 7 October 2026 literature cutoff. The Stata copies have July 2026 PDF
metadata and release-19 example routes, but later retrieval does not authenticate
their exact cutoff-era bytes or an installed licensed build. The older sources
retain their stated manual families and mutable-page limitations. None of these
qualifications licenses a claim that an unevaluated product lacks a feature.

\subsection{Supplied-model computation and the missing causal restriction}
\label{exec22:state-identification}

Let $A\in\{0,1\}$ denote a servicing action and $D\in\{0,1\}$ a subsequent
default. Consider two deliberately small, fully specified mechanisms. In the
selection mechanism, $U\sim\operatorname{Bernoulli}(1/2)$ is an unobserved
favourable credit state, $N\sim\operatorname{Bernoulli}(1/5)$ is independent noise,
and
\begin{equation}
 A=U\mathbin{\oplus}N,\qquad D=1-U.
 \label{exec22:state-selection}
\end{equation}
The exclusive-or operation flips its first binary argument when $N=1$. In the
causal mechanism, an independent action $A\sim\operatorname{Bernoulli}(1/2)$
replaces the first equation, and $D=(1-A)\mathbin{\oplus}N$. Both mechanisms give
exactly the same observable table:
\begin{equation}
\begin{array}{c|cc}
 &D=0&D=1\\ \hline
A=0&1/10&2/5\\
A=1&2/5&1/10
\end{array}
\label{exec22:state-table}
\end{equation}
Consequently, both yield $P(D=1\mid A=0)=4/5$ and
$P(D=1\mid A=1)=1/5$. Under an intervention that replaces only the action
equation, however, the selection mechanism retains default risk $1/2$ for
either action. The causal mechanism has risks $4/5$ and $1/5$. The identified
contrast within each stipulated mechanism is therefore different even though
every observable cell agrees. This is a project derivation by complete finite
enumeration, not a finding about actual customers.~\citep{exec22_state_reference}

With an assumed loss of USD~1,000 per default and action cost USD~100, the net
value of universal action relative to no action is USD~$-100$ under selection
and USD~500 under the causal mechanism. More observations from this same law
cannot distinguish those values. If $w$ is the prior probability of the causal
mechanism in a two-model analysis with these fixed parameters, equal likelihoods
leave its posterior probability equal to $w$ for every observed sample. The
posterior expected margin is
\begin{equation}
 V(w)=600w-100\quad\text{USD},\qquad V(w)>0\ \Longleftrightarrow\ w>1/6.
 \label{exec22:state-prior}
\end{equation}
Priors $w=1/10$ and $w=9/10$ imply USD~$-40$ and USD~440. These numbers expose
the decision's sensitivity to a causal assumption that the likelihood cannot
learn in this example. They do not prohibit Bayesian decisions under declared
assumptions; they require the source and consequences of those assumptions to
remain visible.

BayesiaLab's inspected documentation distinguishes observation from intervention
and names graph surgery and likelihood matching. Its identification discussion
requires causal information expressed in a DAG before estimation. This is a
useful conceptual strength for an adapter. It does not mean that drawing the
action arrow or invoking graph surgery establishes that the bank's graph is
true.~\citep{exec13_ct_bayesia_inference,exec13_ct_bayesia_identification}
An agent could help enumerate the two mechanisms, identify which records might
distinguish them and draft a public experiment. It must retain both until justified
restrictions or new information resolve their difference. A natural-language
statement of confidence cannot replace that missing evidence.

\subsection{State inference must preserve the decision clock}
\label{exec22:state-clock}

The supplied state model may nevertheless be valuable for prediction and decision
support. GeNIe's inspected DBN overview represents a stationary discrete-time
stochastic process. Its discrete-DBN learning page fits parameters by EM in an
already specified structure; the selected page explicitly excludes learning the
dynamic structure. Time-indexed columns must match node identifiers, and existing
network evidence is ignored during that parameter-learning operation. These are
concrete semantic requirements, not interchangeable user-interface details.
\citep{exec13_ct_genie_dbn,exec13_ct_genie_learning}

SMILE's C++ tutorial distinguishes initial and transition conditional-probability
tables, temporal lag order, the number of unfolded slices and integer outcome
indices. A flattened belief array is ordered by slice and outcome. A proposed
Python, Java or R wrapper must preserve those mappings before its output can be
compared with a native model. A matching array shape is insufficient if the
meaning of outcome zero or the initial slice changes.
\citep{exec13_ct_smile_tutorial}

Stata's documented \texttt{sspace} procedure estimates linear state-space
parameters by maximum likelihood using Kalman or diffuse-filter likelihoods.
Its general form separates a latent state $z_t$, measured variables $y_t$,
exogenous inputs $x_t,w_t$ and disturbances:
\begin{equation}
\begin{aligned}
 z_t&=Fz_{t-1}+Bx_t+C\varepsilon_t,\\
 y_t&=Hz_t+Gw_t+J\nu_t.
\end{aligned}
\label{exec22:state-model}
\end{equation}
The coefficient matrices are time invariant in this documented profile.
The stated likelihood uses Gaussian, serially uncorrelated disturbances with
state and observation errors uncorrelated with each other. The documented robust
variance option and stationary/nonstationary initialization methods address
different statistical and numerical choices; they do not certify that a bank
action placed among the regressors is exogenous. A proposed adapter must retain
the declared covariance structure, model constraints and initial-state method.
\citep{exec22_state_stata_sspace}

The postestimation interface makes an especially important distinction:
one-step predictions use past information, filtered states use past and current
information, and smoothed states use all sample information. Its prediction
sample can change the recursive state estimates. The bank's information set
must therefore include availability times, not merely a date column.
\citep{exec22_state_stata_sspace_post}

For an exact example, let a hidden distress state $H$ remain constant for two
dates, with $P(H=1)=1/2$. Conditional on $H$, independent indicators $X_0,X_1$
equal $H$ with probability $3/4$. On observing $X_0=1$, the filtered probability
is $3/4$. If the later observation is $X_1=0$, retrospective smoothing gives
\begin{equation}
 P(H=1\mid X_0=1,X_1=0)
 =\frac{(3/4)(1/4)}{(3/4)(1/4)+(1/4)(3/4)}=1/2.
 \label{exec22:state-smoothing}
\end{equation}
Suppose a separate, stipulated causal model says an action prevents default with
probability $1/10$ when $H=1$, has no effect otherwise, costs USD~60 and avoids
USD~1,000 when successful. Its filtered expected margin is USD~15; using the
later smoothed probability gives USD~$-10$. The latter is a legitimate
retrospective calculation, but it was unavailable when the action was chosen.
The example assumes that the causal benefit law is known; filtering does not
identify that law.~\citep{exec22_state_reference}

Observability is another distinct obligation. In a noise-free two-state model
with $F=(9/10)I$ and $H=(1,1)$, the observability matrix formed by $H$ and $HF$
has rank one. The initial states $(1,0)$ and $(2,-1)$ generate the same output
$(9/10)^t$ at every date, but their first-minus-second component contrasts are
one and three. The null direction $(1,-1)$ remains invisible. Adding a prior
can regularize or select a state representation; it does not make the contrast
observable from these measurements. New measurements, constraints justified by
the application, or a decision defined only on observable combinations are
substantive alternatives. State observability, parameter identification and
causal identification each require their own argument.

\subsection{A posterior, a prediction error and a simulation error answer different questions}
\label{exec22:state-uncertainty}

Stata's \texttt{bayesmh} provides adaptive Metropolis--Hastings estimation with
explicit likelihoods and priors. Its manual cautions that allowed combinations
need not produce a proper posterior, and that no single diagnostic establishes
convergence. Multiple starting points, all model parameters and requested
functions need attention; apparent mixing can conceal multimodal
pseudoconvergence. Chain count, burn-in, adaptation, thinning and random seeds
must therefore be retained in a reproducible specification. The documented
defaults are starting configurations, not universal adequacy guarantees.
\citep{exec22_state_stata_bayesmh}

Netica's inspected learning text instead describes counting, EM and
conjugate-gradient parameter fitting, with prior information supplied through
preexisting tables and experience. Its objective includes both data likelihood
and prior terms, and local optima remain possible. A fitted table should not be
relabeled a sample from a full parameter posterior.~\citep{exec13_latent_netica_algorithms}
The Mplus 8.1 addendum supplies a complementary Bayesian dynamic-SEM profile,
including latent centering and outcome versus residual lag distinctions. Its
missing-covariate rule can remove observations unless the covariate is modeled
as an outcome. A servicing comparison must preserve that modeling decision and
its observation set. The addendum alone does not supply a complete selected
release's prior, sampler and licensing specification.~\citep{exec13_ct_mplus81}

Uncertainty must be named before being compared. For a parameter $\theta$,
training information $\mathcal D$ and a future outcome $Y$, total variance is
\begin{equation}
\begin{aligned}
 \operatorname{Var}(Y\mid\mathcal D)
 &=E[\operatorname{Var}(Y\mid\theta,\mathcal D)\mid\mathcal D]\\
 &\quad+\operatorname{Var}(E[Y\mid\theta,\mathcal D]\mid\mathcal D).
\end{aligned}
\label{exec22:state-variance}
\end{equation}
The first term describes conditional outcome variation; the second propagates
parameter uncertainty. For $\theta\sim\operatorname{Uniform}(0,1)$ and
$Y\mid\theta\sim\operatorname{Bernoulli}(\theta)$ with no data, these terms
are $1/6$ and $1/12$, totaling $1/4$. For two future outcomes conditionally
independent given their shared $\theta$, the variance of their count is
$2(1/6)+4(1/12)=2/3$. Treating the probability as fixed at its posterior mean
$1/2$ gives $1/2$ instead. The shared unknown parameter induces dependence.
\citep{exec22_state_reference}

The selected \texttt{sspace} postestimation manual explicitly states that its
reported prediction RMSE includes disturbance variation but does not account
for parameter-estimation error. It must therefore be labeled accordingly.
An adapter could add a justified parameter-propagation layer, but its result
would be an additional construction requiring its own assumptions and tests.
Neither layer automatically represents graph uncertainty or uncertainty about
the action mechanism.~\citep{exec22_state_stata_sspace_post}

Monte Carlo error is separate again: it measures finite simulation error in a
posterior summary. The \texttt{bayesmh} manual distinguishes effective-sample-size
and batch-means MCSE calculations. Its saved dataset also compresses repeated
states into rows with chain, iteration and frequency fields. Frequencies are
required when summarizing these rows.~\citep{exec22_state_stata_bayesmh}
For values $0.1$ and $0.9$ with frequencies nine and one, the posterior-sample
mean is $0.18$, whereas the unweighted row mean is $0.5$. A CSV conversion that
retains parameter values but drops frequencies silently changes the result.
This exact arithmetic is an export-contract example, not a native execution
test.~\citep{exec22_state_reference}

\subsection{From numerical tolerances to a conditional decision bound}
\label{exec22:state-approximation}

An approximate belief need not support the same action as its reference. Suppose
the exact and approximate probabilities of a beneficial state are $0.61$ and
$0.59$, an action produces USD~100 in that state and costs USD~60. The expected
margins are USD~1 and USD~$-1$. A small probability discrepancy can matter near a
decision threshold even when average predictive scores look similar.

A conditional bound clarifies what an approximation contract would need. Let
$b_t$ and $\widetilde b_t$ be exact and approximate marginal state laws under
the \emph{same fixed policy}. Define total variation as
$\|b-\widetilde b\|_{\rm TV}=\frac12\sum_s|b(s)-\widetilde b(s)|$.
If the reward at each date has range at most $R$, and the marginal discrepancy
is at most $\epsilon_t$, then
\begin{equation}
 |V-\widetilde V|\leq
 R\sum_{t=0}^{T-1}\gamma^t\epsilon_t,\qquad 0\leq\gamma<1.
 \label{exec22:state-bound}
\end{equation}
To see this, subtract the reward minimum, bound each expected-reward difference
by its range times total variation, and apply the triangle inequality to the
discounted sum. With $R=100$ dollars, $\epsilon_t=0.02$, $\gamma=0.9$ and
$T=12$, the bound is
$2(1-0.9^{12})/(1-0.9)$, approximately USD~14.3514.
\citep{exec22_state_reference}

This argument assumes the marginal error bounds. A one-step transition
approximation, a sampler stopping rule or agreement on one benchmark does not
establish them. Nor does a bound for one fixed policy automatically cover the
regret of a different policy selected using approximate beliefs. Deriving
propagation under a specified transition law, verifying it on known-law public
examples, and separating policy-selection error are the next scientific steps.
The proposed product should report these limits instead of attaching an
unqualified confidence label to the action value.

\subsection{Commercial comparison and an executable proposal route}
\label{exec22:state-comparison}

The vendors have different documented strengths. Stata supplies explicit
state/observation equations, initialization options, information-set choices,
Bayesian diagnostics and persistent simulation exports. GeNIe/SMILE supplies
graphical DBNs and concrete native temporal representations. Netica supplies
parameter-learning alternatives with latent or incomplete data. BayesiaLab
distinguishes observation, intervention and identification assumptions. Mplus
supplies a Bayesian dynamic-SEM modeling route. The proposal currently has
exact authored references and semantic requirements; it has no qualified native
adapter for these profiles. Its proposed contribution is to preserve causal
meaning and refusals consistently across the library, CLI and MCP, including
the agent's translation of a business question into a model.

The appropriate reuse route is a conditional, isolated licensed worker with a
versioned exchange contract. A request must identify its task, decision time,
outcome horizon, units, observed and latent variables, action mechanism and
admissible information. A returned result must distinguish filtered from
smoothed states, process from parameter uncertainty, and posterior from Monte
Carlo variation. It must preserve native diagnostics, chain frequencies,
parameter names, initialization and failed calculations. Replacing a vendor's
engine or porting its mathematics is a separate rights and validation decision;
public documentation is not a redistribution entitlement. Prices, commercial
deployment rights and measured operating costs remain unavailable in this
comparison, so no saving or superiority is asserted.

SAS's selected MCMC PDF transfer was incomplete; EViews and HUGIN technical
endpoints failed. Those attempts neither establish their capabilities nor
prove absence. Complete selected-release MCMC/SSM manuals, official state-model
help, HUGIN DBN and intervention documentation, Bayes Server's versioned SDK
contracts, and MATLAB/Wolfram toolbox-specific originals are concrete recovery
targets. Broader complementary searches must also test whether the shortlist
misses a product family. These source obligations can proceed independently of
public prototype design.

The first prototype comparison can use the finite laws above, supplied-state
models with known parameters, and deliberately misspecified alternatives.
Before execution, freeze the query, information clock, causal assumptions,
uncertainty components, monetary inputs and admissible errors. Compare exact
enumeration with a licensed native inference path; test a timestamp violation,
an unobservable contrast, a frequency-loss export, posterior impropriety and a
policy near its decision threshold. Record both correct answers and justified
refusals. Only after such comparisons show merit should the proposal request
costly bank-specific experiments. This route develops useful public evidence
while retaining the unresolved assumptions and commercial terms.
